\bookchapter{פרק 23: חשיבה תכנותית ואוטומציה למתחילים}

\chapterintro{בפרקים הקודמים בנינו לטלפון “חומות הגנה” מאובטחות, וגם התוודענו ל“מוח ה-AI” החכם שלו שבתוך החומות. עכשיו, בואו נגלה את אחד הפוטנציאלים האולטימטיביים של הטלפון: לאפשר לו “לחשוב” כמו מתכנת, ולבצע באופן אוטומטי משימות רציפות שאתה מגדיר. אולי תאמר: “אני לא יודע לתכנת, מה זה קשור אליי?” אל תדאג, הפרק הזה לא דורש ממך לכתוב שורה אחת של קוד – טלפונים מודרניים כבר הורידו את רף התכנות לרמה שניתן להסתדר עם גרירה ושחרור.}

\booksection{מהי בעצם אוטומציה?}

אתה לא צריך לבצע פעולות ידנית בכל פעם, וגם לא צריך לזכור מה לעשות – הטלפון יחליט בעצמו על העיתוי ויבצע בעצמו.\textbf{הנה דוגמה: בכל בוקר כשאתה מתעורר, אתה צריך לכבות ידנית את מצב 'נא לא להפריע', לפתוח את WeChat לבדוק הודעות, לפתוח את אפליקציית מזג האוויר לראות מה ללבוש היום, ולפתוח את NetEase Cloud Music לנגן שיר. אם תארוז את הצעדים האלה לתהליך אוטומציה אחד, ותגדיר “ביצוע אוטומטי בכל בוקר בשעה 7:00”, אז כשאתה מתעורר אתה רק צריך להרים את הטלפון, והכול כבר מוכן. נשמע נהדר, נכון?}אוטומציה היא בעצם שרשרת לוגית של “אם… אז…”, שזה ה-if-else בתכנות. ברגע שתבין את צורת החשיבה הזו, גם אם לעולם לא תכתוב קוד, היכולת שלך לטפל בבעיות בצורה מסודרת תשתפר משמעותית. מרכיבי הליבה:

תנאי הפעלה

: מתי להתחיל? זה יכול להיות “בכל בוקר בשעה 7:00”, “כשאני מתחבר ל-Wi-Fi הביתי”, “כשאני מגיע למיקום מסוים”, “כשסוללת הטלפון יורדת מתחת ל-20\%”.

\begin{enumerate}
  \item \textbf{פעולות לביצוע}: מה לעשות בדיוק? זה יכול להיות לשלוח הודעת WeChat, לפתוח אפליקציה מסוימת, לכוונן עוצמת שמע, להחליף טפט, לנגן מוזיקה, לשלוח התראה… כמעט כל מה שאתה יכול לעשות ידנית בטלפון, אוטומציה יכולה לעשות.
  \item \textbf{לוגיקת ביניים}(אופציונלי אבל עוצמתי): אם מצב A אז לעשות X, אחרת לעשות Y. לדוגמה "אם זה יום עבודה אז לנגן פודקאסט חדשות, בסוף שבוע לנגן מוזיקה", זה נקרא הסתעפות מותנית.
  \item \textbf{2. "קיצורי דרך" ב-iOS – כלי האוטומציה המושלם למשתמשי אפל}אפליקציית "קיצורי דרך" (Shortcuts) של אפל היא אפליקציה מובנית, אין צורך להוריד. הרעיון המרכזי שלה הוא לחבר פעולות שונות יחד כמו אבני בניין. קודם כל, כמה מושגי מפתח:
\end{enumerate}

\booksection{“קיצורי דרך” ב-iOS – כלי האוטומציה המושלם למשתמשי אפל}

: תהליך אוטומציה שלם, שניתן להפעיל בלחיצה ידנית או באמצעות פקודה קולית ל-Siri. לדוגמה, אם תאמר “היי Siri, מצב חזרה הביתה”, היא תבצע אוטומטית את סדרת הפעולות שהגדרת.

\begin{itemize}
  \item \textbf{"אוטומציה"}: פקודה שמופעלת אוטומטית על ידי תנאי הפעלה, ללא צורך בהפעלה ידנית שלך. לדוגמה "כשהטלפון מתחבר ל-Bluetooth של הרכב, פתח אוטומטית את אפליקציית המפות". אוטומציה מוגדרת בלשונית "אוטומציה", נפרדת מ"קיצורי דרך".
  \item \textbf{"פעולה"}: יחידת הביצוע הקטנה ביותר. פתיחת אפליקציה היא פעולה, שליחת הודעה היא פעולה, קבלת מזג אוויר היא פעולה, כוונון בהירות היא פעולה. אתה יכול לשרשר מאות ואלפי פעולות יחד.
  \item \textbf{הנה כמה דוגמאות מעשיות ושימושיות, ברגע שתעקוב אחריהן ותעשה אותן תבין:}דוגמה 1: "שגרת בוקר טוב" בלחיצה אחת
\end{itemize}

בכל בוקר לחיצה אחת (או אמירת “בוקר טוב” ל-Siri), והטלפון יבצע אוטומטית: קבלת מזג האוויר של היום הקראת מזג האוויר והטמפרטורה הצגת לוח השנה של היום פתיחת אפליקציית המוזיקה המועדפת שלך והתחלת ניגון. שלבי הפעולה: קיצור דרך חדש הוסף פעולה “קבל מזג אוויר נוכחי” הוסף “הקרא טקסט” הוסף “קבל אירועי לוח שנה” הוסף “הצג תזכורת” הוסף “פתח אפליקציה” ובחר NetEase Cloud Music. ככה פשוט, הכול בגרירה.

\booksubsection{דוגמה 1: “שגרת בוקר טוב” בלחיצה אחת}

בלשונית “אוטומציה”, צור אוטומציה חדשה בחר “רמת סוללה” הגדר “מתחת ל-30\%” הוסף פעולה “הגדר מצב צריכת חשמל נמוכה לדלוק”, “הנמך בהירות מסך ל-30\%”, “כבה Bluetooth”. מעכשיו כשהסוללה מתחת ל-30\% הטלפון יעבור אוטומטית למצב חיסכון, בלי שתצטרך לדאוג.

\booksubsection{דוגמה 2: “מצב חיסכון בסוללה” אוטומטי}

בכל פעם שאתה משלם בסריקת קוד ב-WeChat/Alipay, הטלפון יקרא את התראת התשלום, וירשום אוטומטית את הסכום ובית העסק בתזכיר או בטבלת Numbers. זה דורש שימוש בתנאי ההפעלה “קבלת הודעת SMS/התראה” באוטומציה של קיצורי דרך (iOS 17+ תומך בקריאת עסקאות ארנק), בשילוב עם הפעולה “הוסף ל-Numbers”. למי שרוצה לעקוב אחר הוצאות אבל מתעצל להקליד ידנית, זה פשוט מושיע.

\booksubsection{דוגמה 3: מעקב הוצאות אוטומטי בסריקת קוד}

3. אוטומציה באנדרואיד – Tasker והחברים שלו

בצד האנדרואיד, המצב קצת יותר מורכב, אבל התקרה גם גבוהה יותר. קודם כל, נתחיל מהפשוט:

\booksection{אוטומציה באנדרואיד – Tasker והחברים שלו}

הרבה טלפוני אנדרואיד מגיעים עם “מקליק אוטומטי” מובנה או שניתן להוריד מחנות האפליקציות. העיקרון סופר-פשוט: אתה מבצע פעולה ידנית פעם אחת, והוא מקליט את הקואורדינטות ומרווחי הזמן של הלחיצות שלך, ואז יכול לשחזר במהירות גבוהה. לדוגמה, אם יש אפליקציה שצריך להיכנס כל יום כדי לקבל פרסי התחברות, תקליט את תהליך הלחיצה פעם אחת, ומעכשיו כל יום זה ירוץ אוטומטית. החיסרון: הוא מזהה רק קואורדינטות, אם פריסת הממשק משתנה (למשל קופצת פרסומת), הוא ילחץ לא נכון.

\booksubsection{רמת מתחילים: מקליק אוטומטי}

Tasker הוא כלי האוטומציה המוכר כחזק ביותר בפלטפורמת אנדרואיד, שולט בראש הטבלה שנים. האמת היא שעקומת הלמידה שלו קצת תלולה, הממשק לא הכי יפה, אבל ברגע שתתחיל, תרגיש שהטלפון פתאום הופך לכל-יכול. גם הליבה של Tasker מבוססת על מודל “הפעלה משימה”:

\booksubsection{רמה מתקדמת: Tasker – מלך האוטומציה של אנדרואיד}

: מגדיר “מתי להפעיל”. תנאי ההפעלה יכולים להיות זמן, מיקום, מצב חיבור Wi-Fi, חיבור Bluetooth, קבלת התראה מסוימת, מצב סוללה, כיוון מסך… תנאי ההפעלה ש-Tasker תומך בהם הם ברמה מטורפת, עשרות סוגים.

\begin{itemize}
  \item \textbf{משימה (Task)}: מגדיר "מה לעשות אחרי ההפעלה". משימה יכולה להכיל מספר פעולות, שמתבצעות לפי סדר. סוגי הפעולות כוללים: כוונון הגדרות מערכת (עוצמת שמע, בהירות, Wi-Fi, Bluetooth, מצב טיסה), שליחת SMS, בקשות HTTP, הרצת פקודות Shell, שליטה בהשמעת מדיה, הצגת תיבת דו-שיח ועוד.
  \item \textbf{הנה כמה תרחישים קלאסיים ושימושיים של Tasker:}הפעלת Wi-Fi אוטומטית כשמגיעים הביתה
\end{itemize}

: הגדר תנאי הפעלה “מחובר לתחנת בסיס מסוימת” (כי ליד הבית אתה מתחבר לתחנת בסיס קבועה) משימה “הפעל Wi-Fi”. לא צריך GPS, לא צורך סוללה, הטלפון מדליק Wi-Fi אוטומטית ברגע שאתה נכנס לאזור השכונה.

\begin{itemize}
  \item \textbf{מצב פגישה}: תנאי הפעלה "כותרת אירוע בלוח שנה מכילה 'פגישה'" משימה "הגדר מצב שקט, כבה עוצמת מדיה, השב אוטומטית ב-WeChat (דורש תוסף)". לעולם לא תדאג יותר שהטלפון יצלצל פתאום בפגישה.
  \item \textbf{הגנה על העיניים בלילה}: תנאי הפעלה "זמן 22:00-06:00" משימה "הפעל מצב הגנת עיניים, הנמך בהירות ל-20\%, הפוך מסך לשחור-לבן (מצב גווני אפור)". למסך שחור-לבן יש השפעה מדהימה על הדחף לגלול בסרטונים קצרים – ברגע שהתמונה הופכת לאפורה, אתה פשוט לא רוצה לגלול יותר.
  \item \textbf{פודקאסט אוטומטי בחיבור אוזניות}: תנאי הפעלה "אוזניות חוטיות מחוברות / אוזניות Bluetooth מחוברות" משימה "פתח אפליקציית XiaoYuzhou והתחל לנגן את הפרק האחרון". מרכיב אוזניות ומיד שומע, בלי לפתוח את הטלפון ולחפש אפליקציה.
  \item \textbf{אם אתה מרגיש ש-Tasker קשה מדי להתחלה, אתה יכול לנסות את האח הפשוט יותר שלו:}MacroDroid
\end{itemize}

. ממשק הרבה יותר ידידותי, תבניות עשירות, מתאים למשתמשים שלא רוצים להתעסק יותר מדי. בנוסף, הרבה מערכות אנדרואיד סיניות (MIUI של Xiaomi, HarmonyOS של Huawei, ColorOS של OPPO) מגיעות עם פונקציות “תרחישים” או “עוזר חכם” מובנות, שלמרות שאינן חזקות כמו Tasker, לגמרי מספיקות לאוטומציה יומיומית.\textbf{MacroDroid}הרווח הגדול ביותר בלימוד אוטומציה הוא לא באמת חיסכון בזמן, אלא

\booksection{יישומים יומיומיים של חשיבה אוטומטית}

. ברגע שתתרגל למודל החשיבה “אם X קורה, אז בצע Y באופן אוטומטי”, תגלה שכמות עצומה של עבודה שחוזרת על עצמה בחיים יכולה להיעלם. לדוגמה:\textbf{בכל פעם שאתה מצלם תמונה, גיבוי אוטומטי לענן ואז דחיסת התמונה המקומית לגודל קטן לחיסכון במקום.}בכל בוקר, קריאה אוטומטית של סידור היום מלוח השנה, יצירת טקסט ושליחתו לבני משפחה דרך WeChat.

\begin{itemize}
  \item כשמתקבלת הודעת SMS על חתימת משלוח, הוספת משימת "לאסוף חבילה" אוטומטית בתזכיר.
  \item כשסוללת הטלפון טעונה במלואה, תזכורת קולית אוטומטית לנתק את המטען (למניעת נזק לסוללה מטעינת יתר).
  \item בכל יום כשמגיע זמן סיום העבודה, אם אתה עדיין ליד המשרד (לפי מיקום תחנת בסיס), תזכורת אוטומטית "הגיע הזמן ללכת הביתה".
  \item כל אלה נשמעים מתקדמים, אבל בעצם ברגע שתשלוט בשני המושגים הבסיסיים של "הפעלה + פעולה", בשילוב עם קיצורי דרך או Tasker, תוכל ליצור את כולם. המפתח הוא
  \item שתחילה תזהה ש"הדבר הזה יכול לעבור אוטומציה"
\end{itemize}

– הרבה אנשים לא שאינם יכולים, אלא שפשוט לא חשבו על זה בכלל.\textbf{5. מדריך להימנעות ממכשולים: בעיות נפוצות באוטומציה}אוטומציה אינה נטולת אתגרים, יש כמה מכשולים שכבר עברנו עבורך:

\booksection{מדריך להימנעות ממכשולים: בעיות נפוצות באוטומציה}

: ההפעלה “הגעה/עזיבת מיקום” ב-iOS משתמשת בטכנולוגיית Geofencing, הדיוק תלוי באות GPS ובאופטימיזציית המערכת. לפעמים הטלפון מופעל רק עשר דקות אחרי שכבר הגעת הביתה. מומלץ להשתמש בהפעלה לפי חיבור Wi-Fi או תחנת בסיס במקום הפעלה לפי מיקום, זה אמין יותר וגם חוסך בסוללה.

\begin{itemize}
  \item \textbf{ה"אוטומציה" של קיצורי דרך דורשת הרשאה}: ב-iOS, רוב ההפעלות בלשונית "אוטומציה" דורשות אישור ידני שלך כדי לרוץ (מנגנון האבטחה של iOS). אבל הפעלה לפי חיבור Wi-Fi, חיבור Bluetooth, ועצירת שעון מעורר – אלה כן יכולים לרוץ אוטומטית לחלוטין, ללא צורך באישור.
  \item \textbf{Tasker דורש הפעלת הרשאת נגישות}: הרבה פונקציות מתקדמות (כמו סימולציית לחיצה) דורשות הרשאת שירות נגישות, מה שיגרום לטלפון להציג התראה קבועה בהגדרות. זהו עיצוב האבטחה של אנדרואיד, אי אפשר להסיר.
  \item \textbf{אל תהפוך את האוטומציה למסובכת מדי}: כשמתחילים עם אוטומציה, תתחיל מתהליכים קטנים של "2-3 פעולות". אל תנסה ישר ליצור תהליך-על של 20 צעדים – ניפוי השגיאות ישגע אותך.
  \item \textbf{טיפ: הגישה הנכונה להתחלת אוטומציה}תקדיש 5 דקות לחשוב: מהן 3 הפעולות שאתה חוזר עליהן הכי הרבה בטלפון כל יום? לדוגמה, פעולות השגרה של כל בוקר, ההרגל אחרי חיבור אוזניות Bluetooth בכל פעם, ההגדרות לפני כל יציאה מהבית… תתחיל מהדברים הקטנים בתדירות גבוהה, ותיצור אוטומציה פשוטה לכל אחד. גם אם כל אוטומציה חוסכת לך רק 5 שניות, לאורך זמן זה מצטבר למשמעותי. וחוץ מזה, תחושת הסיפוק אחרי שיוצרים את זה – באמת ממכרת.
\end{itemize}

\textbf{תרגול: שלושה אתגרי אוטומציה למתחילים}\par

בחר את המתאים לך ביותר והשלם אותו, ברגע שתסיים אתה כבר בפנים:

\begin{tcbpractice}{אתגר 1: אוטומציית בוקר טוב}
: iOS: באפליקציית קיצורי דרך צור קיצור דרך חדש הוסף “קבל מזג אוויר נוכחי” “הקרא טקסט” “פתח אפליקציה” ובחר NetEase Cloud Music. לאחר שתסיים תוכל לומר ל-Siri “בוקר טוב” כדי להפעיל. אנדרואיד: ב-Tasker צור משימה חדשה הוסף “קבל מזג אוויר” “הקרא” “פתח אפליקציה” הגדר זמן הפעלה לכל יום 7:00.

\begin{enumerate}
  \item \textbf{אתגר 2: חיסכון אוטומטי בסוללה נמוכה}: iOS: לשונית אוטומציה חדש "רמת סוללה" מתחת ל-25\% הוסף "הגדר מצב צריכת חשמל נמוכה לדלוק" + "הגדר בהירות 25\%". אנדרואיד: Tasker צור Profile מצב חשמל סוללה 0-25\% משימה "הפעל מצב חיסכון" + "הנמך בהירות".
  \item \textbf{אתגר 3: ניווט אוטומטי בחיבור Bluetooth לרכב}: iOS: אוטומציה חדש "Bluetooth" בחר את Bluetooth הרכב שלך הוסף "פתח אפליקציה" ובחר Gaode Maps + "עוצמת מדיה 80\%". אנדרואיד: Tasker Profile מצב Bluetooth מחובר בחר Bluetooth רכב משימה "פתח Gaode" + "הגבר עוצמת מדיה".
  \item \textbf{נספח לפרק 23: מדריך למתחילים ל-GitHub – מאחסון קוד לפלטפורמת שיתוף פעולה}בפרק הקודם דיברנו על "חשיבה תכנותית", וכדי באמת ליישם חשיבה תכנותית ולעבוד על פרויקטים עם אחרים, אי אפשר להתעלם משם אחד – GitHub. בין אם אתה סטודנט למדעי המחשב שרוצה להשתתף בתחרויות, רוצה לשמור את יצירות הקוד שלך, או פשוט רוצה לראות במה מתעסקים מפתחים מכל העולם, GitHub היא פלטפורמה שאתה חייב להכיר. בחלק זה נתחיל מהמושגים הבסיסיים ביותר, ונבנה הבנה ברורה.
\end{enumerate}

\end{tcbpractice}

\begin{tcbinfo}{מקרה 3: אוטומציה של "שגרת בוקר"}
\begin{itemize}
  \item \textbf{תרחיש}:7:00 תחזית + בטל שקט + חדשות.
  \item \textbf{צעדים}:iOS: Shortcuts > Automation > Time > Weather/Volume/App; Android: MacroDroid > Trigger (7:00) > Volume/URL/TTS.
  \item \textbf{תוצאה צפויה}: מתבצע ללא התקנה, מפחית את רף הכניסה.
\end{itemize}

\end{tcbinfo}
